\documentclass[../../main.tex]{subfiles}% 注意这里的文件路径不能用 ./main.tex ,否则用latexmk编译子文件会报错
\graphicspath{{\subfix{./image/}}} % 指定图片目录,后续可以直接使用图片文件名
% 注意这里的文件路径不能用 ../../image/ ,否则用latexmk编译子文件会报错

% 例如:
% \begin{figure}[H]
% \centering
% \includegraphics[scale=0.3]{图.png}
% \caption{}
% \label{figure:图}
% \end{figure}
% 注意:上述\label{}一定要放在\caption{}之后,否则引用图片序号会只会显示??.

\begin{document}

\section{曲线论基本定理}
曲率和挠率完全决定曲线的形状（不计刚体运动）。
\begin{theorem}[曲线论基本定理]
设$\kappa(s),\tau(s)$为区间$[a,b]$上连续可微函数，且$\kappa(s)>0$，则存在唯一的正则参数曲线$\bm{r}(s)$（不计$\mathbb{E}^3$中的刚体运动），以$s$为弧长参数，$\kappa(s)$为曲率，$\tau(s)$为挠率。
\end{theorem}

\begin{proof}
构造Frenet标架的微分方程组（即Frenet公式），由常微分方程解的存在唯一性定理，标架场唯一确定，其原点轨迹即为所求曲线。
\end{proof}

\begin{corollary}
曲率和挠率均为常数的曲线必为**圆螺旋线**。
\end{corollary}

\begin{example}
给定常曲率$\kappa_0>0$、常挠率$\tau_0$，曲线的参数方程为
\[
\bm{r}(s) = \left( \dfrac{\kappa_0}{\kappa_0^2+\tau_0^2}\cos s, \dfrac{\kappa_0}{\kappa_0^2+\tau_0^2}\sin s, \dfrac{\tau_0 s}{\kappa_0^2+\tau_0^2} \right).
\]
\end{example}




\end{document}